\section{Discussion and Limitations}
\label{sec:discussion}

\textbf{What the detection buys at the application.} The paper's motivation is that a grayhole at $10^{-3}$ to $10^{-4}$ slows every collective that crosses it, and the reader is owed the measurement that tests it. We ran one, and it is unfavourable to the motivation as first stated. We simulated a mixture-of-experts collective at packet level on a 1024-uplink sprayed fabric, with a $10^{-3}$ grayhole injected on a fully loaded uplink, 255 silent drops over the job, and retransmission-timeout recovery at its worst-case setting of 40 ms. The collective's completion time rose by 1.36 percent over the fault-free run (3.588 s to 3.636 s). Because the timeout was at its largest value, the figure should bound every faster recovery setting. Because the simulator's measurement layer has no mitigation actuator, the fault-free run is the only ceiling on what any mitigation could recover: at most 1.36 percent at this operating point. About 99.5 percent of the drops appear to have been absorbed by the collective's pipelining, since a timeout on one chunk overlaps other chunks' slack and only the worst straggler moves the barrier. Consequently, the value of naming a $10^{-3}$ link is not the completion time of the collective that is running when it is named. It lies elsewhere, and we state where without claiming more than the measurement supports. First, a link losing at $10^{-3}$ from a physical cause is a link that is degrading toward $10^{-2}$ or toward a blackhole; Section~\ref{sec:background} already assumes, with CorrOpt~\cite{corropt}, that such causes persist until repaired, so early naming has repair and attribution value that a single completion-time measurement cannot capture. Second, the 1.36 percent is on one collective, and a fleet runs thousands of links continuously, so the aggregate is not 1.36 percent of one job. Last, the contribution of this paper is the measurement instrument and its boundary. The application-level value of low-loss detection is an open question, and our own simulation bounds it at about 1.4 percent per collective under timeout-dominated recovery.

\textbf{Common-mode and congestion loss.} The failures of Section~\ref{sec:robustness} are the ledger's largest limitation, and they are properties of the decision rule, not of the witness. The registers still recover each link's loss exactly under a fleet-wide shift and under incast. What fails is a comparison against a floor that moves with the fault in the first case, and the absence of any question about the cause of the loss in the second. A fix for the first may need a guard that detects when the floor itself is drifting, or a corroboration gate that requires a link's ratio to stand out relative to its neighbours as well as relative to the floor, i.e., the relative normalization that lets SprayCheck hold in that regime. A fix for the second needs a signal that separates queue loss from link loss; the testbed exposes per-packet queue depth in the egress pipeline, so a stamp that carries it, or a gate that consults it, is the obvious candidate. Both are design work with their own evaluation, and this paper does not do them.

\textbf{What the two bytes buy, and what they do not.} Equation~\eqref{eq:loss} contains no information a zero-byte counter pair lacks, and Section~\ref{sec:detection} shows the two are identical when the pair is read at one instant. What the stamp buys is that the sender's count arrives in band, so the receiver holds both terms of the difference and reads them in one operation. The pair's tolerance to read timing is about ten microseconds across two switches, against a measured 2.6 ms for the nearest pairwise read and 350 ms for a census on our controller. The stamp does not buy a better inference: the localization score of 1.00 in Section~\ref{sec:localization} is definitional given per-directed-link counts, and we present the ledger there as the oracle row against which the passive detectors' aliasing is measured, not as a result.

\textbf{A synchronized counter snapshot.} The read-skew argument invites an obvious reply: do not read the two registers at one instant, latch them at one instant and read them at leisure. If both switches copied their per-directed-link counters into shadow registers on a common epoch boundary derived from precision time, the controller could take its full 350 ms to collect them. The skew term would then be sub-microsecond, an order below the tolerance Section~\ref{sec:detection} establishes, and the counter pair would match the witness at zero wire bytes. This is the alternate-marking family of loss measurement~\cite{rfc9341}, and RFC 6374 is built around it. We believe the reply holds in principle, and we scope the paper against it. Three things bound it in practice. First, the snapshot is not free in the data plane, and our own history prices it. The program this ledger replaced kept alternating counter banks per epoch with a guard interval. The compile gate that retired it measured the banked scheme at two more tables, three more SRAM blocks, one more stateful ALU, and two more header containers than the ledger at the same stage count, before any time synchronization. It also carried the bank parity in a fabric header byte, so it was not zero-byte either, and a transit stage that failed to preserve that byte produced a 56 percent false-blackhole rate until it was found. A snapshot that keys on time rather than on a carried bit avoids the byte but not the banks. Second, the witness needs no fabric-wide coordination: no time distribution, no trusted common time base, and no agreement between two switches about when an epoch began. That is the operational asymmetry that survives the reply. Third, and stated plainly, an operator who has precision time on every switch and can afford the second register bank and the bank swap should prefer the synchronized counter pair, and the witness is then unnecessary. We did not build the synchronized variant; its stage cost on this program is an estimate from the banked scheme we did build.

\textbf{The wire reduction, reconsidered.} Halving the header from four bytes to two cost the program's last ingress stage (Section~\ref{sec:cost}), and the soak below shows it also introduced an anomaly the four-byte program did not have and a structural aliasing of traffic that carries no real spray choice. Neither header escapes the 16-bit wrap bound, since both carry a 16-bit sequence. An operator who weights pipeline headroom or the open anomaly over one byte per packet per hop can reasonably prefer the four-byte program, and the evidence in this paper now points that way rather than merely permitting it. We report both because the trade should be made rather than assumed.

\textbf{The anomaly on unmeasured sublinks, as far as the evidence goes.} During a 57-cycle soak of the two-byte program, the injected-loss recovery passed in every cycle. However, two sublinks that were neither armed nor measured each showed one stamped packet with no matching arrival during clean traffic, a signature not seen across about 3\,200 historical cycles of the four-byte program. We took the investigation as far as the shared hardware allowed. Physical-layer loss was ruled out by port-level media-access counters (9\,224 frames sent and received) and ambient noise by a fully idle window that moved no register. A byte-identical build of the four-byte source, deployed through the same pipeline and run with the same script immediately afterward, showed 0 anomalies in 100 cycles against 2 in 57 for the two-byte build, which isolates the anomaly to the wire-reduction build rather than the testbed. A probe with no injector table writes reproduced it, which rules out the control-plane writes that accompanied the first two occurrences. An idle-then-burst trigger looked confirmed at 3 of 5 post-idle bursts, then two further runs with the corrected instrument found 0 of 9. The honest count is therefore 3 of 14, and the surviving hypothesis is a genuinely cold fabric after a fresh program load rather than mid-session idleness. The root cause is not established, and the fresh-load protocol that would settle it needs its own bring-up budget. Every number in this paper comes from armed, measured sublinks with baseline and post-traffic reads, and none is affected. However, a claim of zero false positives must be read beside a one-packet phantom loss that recurs on the same program, at a rate near 2 per 57 cycles, on unmeasured sublinks.

\textbf{A structural aliasing introduced by the wire reduction.} While restoring the two-byte program after the comparison above, a census of traffic classes outside the probe's set showed phantom arrivals with no stamps on two classes, reproduced across two independent bring-ups. The explanation is mechanistic. Bring-up's own port-verification frames carry no real spray choice and arrive with a spray field of zero, and the new ingress reconstruction maps every frame with spray zero to the first virtual link regardless of which leaf sent it, where the old design read the link identifier the sending egress had written. On the testbed this never touched a measured sublink. In a production fabric, control and management traffic that is stamped but not sprayed would be attributed to one link's counters, and the two-byte design needs either a table entry that excludes unsprayed traffic from the stamp or a spray value reserved for it. We record this as a correctness cost of the reduction that Section~\ref{sec:cost} did not price.

\textbf{Silicon scope.} Section~\ref{sec:silicon} states the limits of the silicon cells: a virtual fabric, one percent of a port, one sublink, computed rather than measured added load, and a one-packet discrepancy explained but not re-measured. The stage budget is that of a stripped program, forwarding plus the ledger. A production leaf program that also carries load balancing, transport offload, access control, tunnelling, and quality of service does not have a spare stage on a Tofino~1, and whether the ledger fits beside such a program, or on a Tofino~2, is not measured here.

\textbf{Generality.} The replay fabric has four leaves and eight spines with uniform spray. FlowPulse's dilution grows with senders per port and SprayCheck's intersection improves with leaves, so we expect the ordering of the arms to survive a change of radix while the crossover rates are properties of this fabric. The $1/p^2$ form of Proposition~\ref{prop:cost} does not depend on the topology; its constants do, through the floor and the epoch volume, and Section~\ref{sec:detection} sweeps both.
